\documentclass[12pt,a4paper]{article}

% ==================================================
% 1. CONFIGURATION
% ==================================================

% Enable / disable title page
\newif\ifshowtitlepage
\showtitlepagetrue
% \showtitlepagefalse

% Enable / disable table of contents
\newif\ifshowtoc
\showtocfalse

% Enable / disable bibliography
\newif\ifshowbibliography
\showbibliographyfalse

% Document information
\newcommand{\worktitle}{Backlog Management Methodology}
\newcommand{\worksubtitle}{NAO Robot Exercise System}
\newcommand{\teamname}{NAOMI}
\newcommand{\subject}{NAO Robot Exercise System}

\newcommand{\university}{Slovak University of Technology in Bratislava}
\newcommand{\faculty}{Faculty of Informatics and Information Technologies}

\newcommand{\academicyear}{2026/2027}
\newcommand{\location}{Bratislava}
\newcommand{\documentyear}{2026/2027}

% Team leaders
\newcommand{\leaderone}{Ing. Jakub Perdek, PhD.}
\newcommand{\leadertwo}{doc. Ing. Fedor Lehocki, PhD.}

% Team members
\newcommand{\memberone}{Bc. Akhmadinurov Amal}
\newcommand{\membertwo}{Bc. Danilov Denis}
\newcommand{\memberthree}{Bc. Dmitriev Anton}
\newcommand{\memberfour}{Bc. Gainetdinov Ruslan}
\newcommand{\memberfive}{Bc. Kuzin Timofei}
\newcommand{\membersix}{Bc. Tikhoblazhenko Alena}
\newcommand{\memberseven}{Bc. Zimnitski Aliaksei}


% ==================================================
% 2. PACKAGES
% ==================================================

\usepackage[
    a4paper,
    top=3cm,
    bottom=2.7cm,
    left=2.8cm,
    right=2.8cm,
    headheight=45pt,
    headsep=20pt
]{geometry}

\usepackage[T1]{fontenc}
\usepackage[utf8]{inputenc}
\usepackage[english]{babel}
\usepackage{lmodern}

\usepackage{graphicx}
\usepackage[table]{xcolor}
\usepackage{fancyhdr}
\usepackage{titlesec}
\usepackage{setspace}
\usepackage{lastpage}
\usepackage{array}
\usepackage{float}
\usepackage{caption}
\usepackage{booktabs}
\usepackage{tabularx}
\usepackage{lipsum}
\usepackage{csquotes}

% Bibliography
\usepackage[
    backend=biber,
    style=ieee,
    sorting=none
]{biblatex}

\addbibresource{references.bib}

% Hyperlinks
\usepackage[hidelinks]{hyperref}


% ==================================================
% 3. COLORS
% ==================================================

\definecolor{primary}{HTML}{000000}
\definecolor{secondary}{HTML}{000000}
\definecolor{muted}{HTML}{000000}
\definecolor{linecolor}{HTML}{000000}
\definecolor{naomiTableHeader}{HTML}{17445A}
\definecolor{naomiTableStripe}{HTML}{EDF7FA}
\definecolor{naomiTableRule}{HTML}{66BBD0}

% Ragged-right columns keep narrow table cells readable.
\newcolumntype{L}[1]{>{\raggedright\arraybackslash}p{#1}}
\newcolumntype{Y}{>{\raggedright\arraybackslash}X}


% ==================================================
% 4. TYPOGRAPHY
% ==================================================

\setstretch{1.2}

\setlength{\parindent}{0pt}
\setlength{\parskip}{7pt}

% Section
\titleformat{\section}
    {\Large\bfseries\color{primary}}
    {\thesection}
    {0.8em}
    {}

% Subsection
\titleformat{\subsection}
    {\large\bfseries\color{primary}}
    {\thesubsection}
    {0.8em}
    {}

% Subsubsection
\titleformat{\subsubsection}
    {\normalsize\bfseries\color{primary}}
    {\thesubsubsection}
    {0.8em}
    {}

\titlespacing*{\section}{0pt}{20pt}{10pt}
\titlespacing*{\subsection}{0pt}{14pt}{7pt}
\titlespacing*{\subsubsection}{0pt}{12pt}{5pt}

% Captions
\captionsetup{
    font=small,
    labelfont=bf,
    justification=centering
}

% Bibliography title
\DefineBibliographyStrings{english}{
    bibliography = {References}
}


% ==================================================
% 5. HEADER AND FOOTER
% ==================================================

\pagestyle{fancy}
\fancyhf{}

% Header logo
\fancyhead[L]{
    \includegraphics[
        width=4cm,
        height=1.15cm,
        keepaspectratio
    ]{assets/horizontal.png}
}

% Header information
\fancyhead[R]{
    \begin{tabular}{r}
        {\footnotesize\bfseries \worksubtitle}\\[3pt]
        {\scriptsize Academic Year \academicyear}
    \end{tabular}
}

% Header line
\renewcommand{\headrulewidth}{0.7pt}

% Footer left
\fancyfoot[L]{
    {\scriptsize
    \teamname\enspace|\enspace\subject}
}

% Footer right
\fancyfoot[R]{
    {\scriptsize
    \thepage\enspace/\enspace\pageref{LastPage}}
}

% Footer line
\renewcommand{\footrulewidth}{0.4pt}

% Plain pages
\fancypagestyle{plain}{
    \fancyhf{}

    \fancyhead[L]{
        \includegraphics[
            width=4cm,
            height=1.15cm,
            keepaspectratio
        ]{assets/horizontal.png}
    }

    \fancyhead[R]{
        \begin{tabular}{r}
            {\footnotesize\bfseries \worksubtitle}\\[3pt]
            {\scriptsize Academic Year \academicyear}
        \end{tabular}
    }

    \fancyfoot[L]{
        {\scriptsize
        \teamname\enspace|\enspace\subject}
    }

    \fancyfoot[R]{
        {\scriptsize
        \thepage\enspace/\enspace\pageref{LastPage}}
    }

    \renewcommand{\headrulewidth}{0.7pt}
    \renewcommand{\footrulewidth}{0.4pt}
}




% ==================================================
% 6. TITLE PAGE
% ==================================================

\newcommand{\makemytitlepage}{
    \begin{titlepage}
        \thispagestyle{empty}
        \centering

        % ==========================================
        % UNIVERSITY LOGO
        % ==========================================

        \vspace*{0.3cm}

        \includegraphics[
            width=6cm,
            height=5cm,
            keepaspectratio
        ]{assets/vertical.png}

        \vspace{0.5cm}

        {\large\bfseries
        \university\par}

        \vspace{0.2cm}

        {\normalsize
        \faculty\par}

        % ==========================================
        % PROJECT TITLE
        % ==========================================

        \vspace{1.7cm}

        {\small\MakeUppercase{\worksubtitle}\par}

        \vspace{0.5cm}

        \begin{minipage}{0.95\textwidth}
            \centering

            {\fontsize{25}{32}\selectfont\bfseries
            \worktitle\par}
        \end{minipage}

        % ==========================================
        % TEAM INFORMATION
        % ==========================================

        \vfill

        \begin{minipage}{0.9\textwidth}
            \centering
            \begin{minipage}[t]{0.46\linewidth}
                \centering
                {\normalsize\bfseries Team Leaders\par}
                \vspace{0.25cm}
                {\small
                \leaderone\par
                \leadertwo\par}
            \end{minipage}%
            \hfill
            \begin{minipage}[t]{0.48\linewidth}
                \centering
                {\normalsize\bfseries Team Members\par}
                \vspace{0.25cm}
                {\small
                \memberone\par
                \membertwo\par
                \memberthree\par
                \memberfour\par
                \memberfive\par
                \membersix\par
                \memberseven\par}
            \end{minipage}
        \end{minipage}

        % ==========================================
        % BOTTOM INFORMATION
        % ==========================================

        \vfill

        {\small
        \location\enspace\textbullet\enspace\academicyear\par}

    \end{titlepage}
}

% ==================================================
% 7. IMAGE PLACEHOLDER
% ==================================================

% Usage:
% \imageplaceholder{Text}{Caption}

\newcommand{\imageplaceholder}[2]{

    \begin{figure}[H]
        \centering

        \begingroup
        \setlength{\fboxsep}{0pt}

        \fbox{
            \begin{minipage}[c][5cm][c]{0.85\textwidth}

                \centering

                {\Large\bfseries IMAGE}

                \vspace{0.3cm}

                {\normalsize #1}

            \end{minipage}
        }

        \endgroup

        \caption{#2}

    \end{figure}
}


% ==================================================
% 8. REAL IMAGE COMMAND
% ==================================================

% Usage:
% \documentimage{path}{caption}{label}

\newcommand{\documentimage}[3]{

    \begin{figure}[htbp]
        \centering

        \includegraphics[
            width=0.85\textwidth,
            keepaspectratio
        ]{#1}

        \caption{#2}
        \label{#3}

    \end{figure}
}


% ==================================================
% 9. TABLES
% ==================================================

% Usage:
% \begin{naomitable}{Caption}{label}
%     \begin{tabularx}{\textwidth}{L{2.8cm} Y L{2.6cm}}
%         \rowcolor{naomiTableHeader}
%         \textcolor{white}{\bfseries Name} & ... \\
%         ...
%         \arrayrulecolor{naomiTableRule}\bottomrule[1pt]
%     \end{tabularx}
% \end{naomitable}

\newenvironment{naomitable}[2]{%
    \begin{table}[htbp]
        \centering
        \caption{#1}
        \label{#2}
        \begingroup
        \small
        \renewcommand{\arraystretch}{1.4}
        \setlength{\tabcolsep}{7pt}
        \rowcolors{2}{naomiTableStripe}{white}
}{%
        \endgroup
    \end{table}
}

% Usage:
% \tableplaceholder{Caption}{label}

\newcommand{\tableplaceholder}[2]{
    \begin{naomitable}{#1}{#2}
        \begin{tabularx}{\textwidth}{L{2.8cm} Y L{2.6cm}}
            \rowcolor{naomiTableHeader}
            \textcolor{white}{\bfseries Name} &
            \textcolor{white}{\bfseries Description} &
            \textcolor{white}{\bfseries Status} \\
            Item A & Sample description & Completed \\
            Item B & Sample description & In Progress \\
            Item C & Sample description & Planned \\
            \arrayrulecolor{naomiTableRule}\bottomrule[1pt]
        \end{tabularx}
    \end{naomitable}
}


% ==================================================
% 10. DOCUMENT
% ==================================================

\begin{document}

\ifshowtitlepage
    \makemytitlepage
\fi

\ifshowtoc
    \tableofcontents
    \clearpage
\fi

% ==================================================
% 1. PURPOSE
% ==================================================

\section{Purpose of the Document}

This document describes how the NAOMI team manages its project backlog
using Jira and a Scrumban approach. The aim is to make planned work
visible, keep responsibilities clear, and support regular delivery of
results throughout the academic year.

Scrumban combines sprint-based planning and review with the visual
workflow of a Kanban board. For NAOMI, this means agreeing on short-term
goals while allowing the team to monitor work in progress, identify
blockers, and adjust priorities when necessary.

% ==================================================
% 2. SCOPE
% ==================================================

\section{Scope and Responsibilities}

The methodology applies to all NAOMI team members working on analysis,
robot-control software, motion processing, testing, documentation, and
other project activities. Jira serves as the central place for recording
and following the status of these activities.

Every member is expected to keep their assigned issues up to date.
The team jointly reviews priorities, dependencies, and progress.
Coordination responsibilities are assigned to a team leader.

% ==================================================
% 3. BACKLOG STRUCTURE
% ==================================================

\section{Backlog Organisation in Jira}

\subsection{Product Backlog}

The product backlog contains the features, improvements, research tasks,
experiments, and technical work that may be needed to reach the project's
goals. Items should have a concise title and a description of the expected
outcome. Larger activities may be represented as Epics and divided into
smaller Stories or Tasks, using the issue types available in the team's
Jira project.

The backlog is reviewed and reordered when requirements change or new
information becomes available. Higher-priority items should be
sufficiently clear before the team starts working on them.

\subsection{Sprint Backlog}

At the beginning of each sprint, the team selects a realistic set of
issues from the product backlog, taking into account available capacity,
priorities, and dependencies. The selected work is visible on the Jira
board. Unfinished issues are discussed during the sprint review or next
planning session and are not automatically treated as completed.


% ==================================================
% 4. SCRUMBAN WORKFLOW
% ==================================================

\section{Scrumban Workflow}

The NAOMI team uses a Jira board to manage tasks and
track their progress throughout the development process.
Our workflow consists of five stages, allowing team
members to monitor the current state of each issue.

\begin{naomitable}{Scrumban workflow states}{tab:workflow}
    \begin{tabularx}{\textwidth}{L{2.8cm} Y}
        \rowcolor{naomiTableHeader}
        \textcolor{white}{\bfseries Status} &
        \textcolor{white}{\bfseries Description} \\
        To Do & The task is planned but work has not started. \\
        In Progress & A team member is actively working on the task. \\
        In Review & The implementation is completed and waiting
        for review or feedback from other team members. \\
        Testing & The task has passed the review stage and is being tested
        to verify that it works correctly and meets the requirements. \\
        Done & The task has been reviewed, tested, and successfully completed. \\
        \arrayrulecolor{naomiTableRule}\bottomrule[1pt]
    \end{tabularx}
\end{naomitable}

Team members are responsible for keeping the status
of their assigned issues up to date in Jira.
Tasks should be moved between columns as soon as
their status changes, rather than only at the end
of a sprint.

If a task cannot be completed due to a blocker,
the responsible team member should document the
problem in Jira and inform the rest of the team.

Following the Scrumban approach, the team focuses
on completing tasks already in progress before
starting new ones. This helps avoid unnecessary
multitasking and improves the overall workflow.


% ==================================================
% 5. REGULAR REVIEW
% ==================================================

\section{Reviewing Progress and Identifying Blockers}

The Jira board is checked during regular team meetings and whenever an
assignee makes meaningful progress. During a review, the team considers
whether the current sprint goal is still achievable and whether work is
moving through the board as expected.

\begin{itemize}
    \item \textbf{Priorities:} Are we working on the most important issues?
    Have new requirements changed the plan?
    \item \textbf{Progress:} Which issues have remained In Progress for too
    long? Is anything blocked or waiting for a decision?
    \item \textbf{Ownership:} Does each active issue have a clear assignee?
    Are issue descriptions and next steps understandable?
    \item \textbf{Delivery:} Are completed results tested, documented, and
    ready to demonstrate or integrate?
\end{itemize}

Issues that need support should be raised at the next team meeting or
communicated earlier through the team's agreed communication channel.
Critical blockers should not be left unresolved until the sprint ends.

% ==================================================
% 6. ISSUE RULES
% ==================================================

\section{Rules for Managing Jira Issues}

\subsection{Creating and Updating Issues}

An issue should describe a specific activity or expected result. It
should normally include a title, a short description, an assignee when
assigned, and relevant acceptance criteria. Links to code, designs,
experiments, or documentation may be attached where useful.

Each member updates the status of their issues and records important
findings or blockers in Jira. If an issue is too broad or cannot be
completed within a reasonable period, it should be divided into smaller
pieces of work.

\subsection{Changes During a Sprint}

Scrumban allows the team to react to new information without losing
visibility of its commitments. Urgent work may be added during a sprint
after discussion with the team, particularly if it affects the sprint
goal or another member's progress. The reason for a significant change
should be documented in the related Jira issue.

\subsection{Finishing Work}

An issue is moved to Done only when its agreed acceptance criteria are
satisfied. Depending on the type of work, this may involve a code review,
a successful test, an updated document, or a demonstration of the result.
The team should use its Definition of Done, once agreed, as the common
reference for deciding whether an item is finished.

% ==================================================
% 7. PLANNING AND ESTIMATION
% ==================================================

\section{Planning and Estimation}

The team estimates work when an estimate is useful for sprint planning
and prioritisation. Story points may be used to express relative effort,
complexity, and uncertainty. An issue type does not determine a fixed
number of story points: for example, different Tasks or Stories can
require very different amounts of work.

\begin{naomitable}{Illustrative relative effort scale for discussion}{tab:estimates}
    \begin{tabularx}{\textwidth}{L{2.4cm} Y}
        \rowcolor{naomiTableHeader}
        \textcolor{white}{\bfseries Points} &
        \textcolor{white}{\bfseries Interpretation} \\
        1 & Very small, well-understood task \\
        2 & Small task with limited uncertainty \\
        3 & Moderate task requiring several steps \\
        5 & Larger task involving more effort or uncertainty \\
        8 & Complex activity that may need to be split \\
        \arrayrulecolor{naomiTableRule}\bottomrule[1pt]
    \end{tabularx}
\end{naomitable}

This scale is a \textbf{proposal}, not an already approved NAOMI rule.
The team can adopt or change it after discussing a few example issues.
Capacity estimates should also account for university deadlines and the
availability of individual members. At the end of a sprint, completed
work, carried-over tasks, and blockers can be used to improve the next
planning session.

% ==================================================
% 8. CONTINUOUS IMPROVEMENT
% ==================================================

\section{Continuous Improvement}

At the end of each sprint, the team briefly reviews what worked well,
what slowed progress, and which changes would help in the next sprint.
Possible improvements include clearer Jira issue descriptions, more
consistent board updates, smaller tasks, and adjustments to meeting or
review practices.

The purpose of this methodology is not to introduce unnecessary
administration. Jira should give the team a reliable overview of its
work while leaving enough time for actual development and experimentation.

% ==================================================
% OPTIONAL REFERENCES
% ==================================================

\ifshowbibliography
    \clearpage
    \printbibliography[
        heading=bibintoc,
        title={References}
    ]
\fi

\end{document}
